\chapter{排除犯罪事由I：排除违法事由}

% Source PDF physical page 88; see audit/source_coverage.json.
\begin{EditionOpeningPanel}
1.行为人具备了客观要件和主观要件后，就意味着制造了法益侵害事实，简称违法事实，其行为便具有法益侵害性，简称违法性。但这只是暂时结论。如果行为人具有排除违法事由，则意味着行为被正当化了。排除违法事由也被称为违法阻却事由。

2.考点重要级别：正当防卫\ensuremath{>}紧急避险\ensuremath{>}被害人承诺。
\end{EditionOpeningPanel}

\section{正当防卫}

[问题]小芳的父亲晚上看到窗户外面有人影，以为是小偷，将对方打倒在地，导致轻伤。实际上对方是小芳的男友狗蛋，前来与小芳幽会。对小芳的父亲该如何处理？（幽会被打案）

第20条第1款为了使国家、公共利益、本人或者他人的人身、财产和其他权利免受正在进行的不法侵害，而采取的\CJKunderwave{制止不法侵害}的行为，对不法侵害人造成损害的，属于正当防卫，不负刑事责任。

\InsertDiagram{figures/manual_table_p088_1.tex}

[注意]体系地位（四星级考点）。正当防卫是排除违法事由，那么就需要事先存在违法性（法益侵害性），然后再通过“正当防卫”将违法性排除掉。先有“尘土”，再将“尘土”扫掉。因此，讨论一个行为是不是正当防卫，前提是该行为已经符合客观要件，具有违法性（法益侵害性），具有“犯罪行为”的表面特征。如果一个行为不符合客观要件，就没必要讨论该行为是不是正当防卫，即使该行为也起到制止不法侵害的效果。亦即，审查步骤走不到“排除违法事由”这个板块。

% Source PDF physical page 89; see audit/source_coverage.json.
例如，甲入室盗窃，主人乙发现后喊了一声：“谁？”甲便逃离。乙的喊声虽然制止了不法侵害，但其行为本身不是危害行为，没有法益侵害性，就没必要讨论其是不是正当防卫。换言之，没必要用正当防卫来论证该行为是无罪的。

[巩固练习］题1，甲向乙扔飞刀，乙躲闪开。乙的躲闪行为不是危害行为，没有法益侵害性。因此，不需要用正当防卫来论证该行为是无罪的，亦即该行为不是正当防卫。

题2，甲抢劫了乙的钱包后，骑摩托车逃跑。乙为了夺回钱包而开车追赶，甲骑摩托车撞到护栏，死亡。乙的追赶行为本身是合法正当的，本身没有法益侵害性，在客观要件中便可得出无罪结论，不需要讨论是不是正当防卫（2012年第7题）。如果乙开车撞倒甲，致甲重伤，夺回钱包，则乙的行为初步来说具有法益侵害性，具有“抢劫犯罪行为”的表面特征，警方会以为乙在抢劫，接下来需要讨论乙是不是正当防卫。在此种情况下，乙是正当防卫（2007年第2题）。

\subsection*{一、起因条件}

正当防卫面临的侵害具有不法性、客观性、现实性。

\subsubsection*{（一）不法性}

1.不法侵害行为仅限于人的行为，因为只能针对人的行为进行合法和不法的评价。山洪暴发、地震灾害、野狗咬人等只能视为单纯的危险，只能进行紧急避险。

2.防卫人不限于被害人本人。只要是面临不法侵害，不管是被害人本人，还是无关第三人，都可以正当防卫，予以制止。例如，甲在路上看到乙在强奸妇女，甲对乙可以正当防卫。

3.不法侵害包括对国家法益、公共法益的不法侵害。例如，外国间谍甲正在窃取我国国家秘密，群众乙看到，情况紧急，乙对甲可以反击防卫。（2024年试题）

4.事态具有紧急性。正当防卫是私力救济。原则上，只能在来不及寻求公力救济（司法程序的救济）的情况下才允许私力救济。所以，正当防卫制止的不法侵害，要具有来不及寻求公力救济的紧急特征。例如，甲看到官员在贪污，甲不能对其进行防卫反击，因为对这种不法侵害，完全可以寻求公力救济，例如向司法机关检举。

\subsubsection*{（二）客观性}

不法侵害，只要求是客观的不法侵害即可。至于在“不法侵害”前加什么修饰语并不重要。

1.客观要件的修饰语

无论是对作为的或不作为的不法侵害进行制止，都构成正当防卫。例如（2019年试题），甲开车不慎撞伤乙，欲逃逸，行人丙强迫甲救助乙，属于针对不作为的不法侵害的正当防卫。

2.主观要件及排除责任的修饰语

无论是故意的或过失的不法侵害，还是对无责任年龄者的或无责任能力者的不法侵害进行制止，都构成正当防卫。

% Source PDF physical page 90; see audit/source_coverage.json.
\InsertDiagram{figures/manual_table_p090_1.tex}

例如，聋哑人甲在狩猎时，误将乙当作野兽准备射击，站在甲身后较远的丙发现甲的行为，赶紧开枪将甲打伤，保护了乙的生命。甲的行为属于过失的不法侵害，对其可以进行正当防卫。（2009年第3题）

[注意］面对精神病患者、未成年人的不法侵害，多数说认为（司法解释），能够正当防卫；\footnote{2020年8月28日最高人民法院、最高人民检察院、公安部《关于依法适用正当防卫制度的指导意见》规定：“成年人对于未成年人正在实施的针对其他未成年人的不法侵害，应当劝阻、制止；劝阻、制止无效的，可以实行防卫”，“明知侵害人是无刑事责任能力人或者限制刑事责任能力人的，应当尽量使用其他方式避免或者制止侵害；没有其他方式可以避免、制止不法侵害，或者不法侵害严重危及人身安全的，可以进行反击”。这些规定表明，司法解释认为无责任年龄人、无责任能力人的侵害行为是不法侵害。同时，刑法第20条表明，正当防卫行为就是制止不法侵害的行为。所以，司法解释倾向于认为，对无责任年龄人、无责任能力人的不法侵害，可以进行正当防卫。}少数说认为，只能紧急避险。（2014年第52题）

[提示】原则上，面临不法侵害，防卫人没有躲避义务，也即“正义不需要向不正义让步”。但是，例外的是，遇到精神病患者、未成年人实施的不法侵害，如果不是严重危及人身安全的不法侵害，应先采取躲避措施；若无法躲避，则可以反击防卫。如果是严重危及人身安全的不法侵害，则不需先躲避，可以直接防卫反击（司法解释规定）。

[总结]动物侵害问题：（2009年第3题）

第一，有主人的狗，主人唆使狗咬人，属于主人故意的不法侵害。反击狗，属于正当防卫。

第二，有主人的狗，主人由于管理过失，导致狗跑出来自发咬人，这属于主人过失的不法侵害。反击狗，属于正当防卫。

第三，有主人的狗，主人没有管理过失，没有任何疏忽，由于不可抗力（如发生地震）导致狗跑出来自发咬人。这种情形不属于主人的不法侵害，主人没有制造法律不允许的危险。反击狗，属于紧急避险。

第四，无主人的野狗自发咬人。由于合法与不法的评价只能针对人的行为，对这种狗咬人的行为不能认为是不法侵害，只能认为是一种单纯的危险。反击狗，属于（防御型）紧急避险。\footnote{有观点认为，由于野生动物不存在需要保护的法益，不存在保护较大法益、损害较小法益的问题，所以反击它不是紧急避险。其实不能一概而论。一般的野生动物是生态环境的一部分，具有值得保护的法益，除非对人类完全有害。例如，蚊子咬人，拍死它，不算紧急避险，因为蚊子是害虫，其本身不存在需要保护的法益。}

\subsubsection*{（三）现实性}

不法侵害必须是现实存在的。如果不存在现实的不法侵害，行为人误以为存在不法侵害，并进行所谓的防卫，就是假想防卫。假想防卫的特点是，\CJKunderwave{好心办坏事}。

例1，在本节开头“幽会被打案”中，客观上不存在不法侵害，因此小芳父亲属于假想防卫。

% Source PDF physical page 91; see audit/source_coverage.json.
例2，甲看到乙在不法侵害丙，便上前将乙扑倒，实际上乙是便衣警察，在抓捕丙。甲属于假想防卫。

例3，小芳与情夫狗剩正在家中苟且，丈夫狗蛋回家，小芳情急之下指着狗剩喊：“他强奸我，赶紧抓！”狗蛋不知情，将狗剩打伤。狗蛋属于假想防卫。小芳属于故意伤害罪的间接正犯。

对于假想防卫的处理：第一，假想防卫不可能是故意为之，否则就是故意犯罪而非假想防卫。第二，假想防卫，如果防卫人有过失，就是过失犯罪。第三，假想防卫，如果防卫人没有过失，就只能按照意外事件处理。

\subsection*{二、时间条件}

正当防卫制止的是正在进行的不法侵害。如果不法侵害尚未开始，便对其防卫，属于事前防卫；如果不法侵害已经结束，对其防卫，属于事后防卫。二者都属于防卫不适时，不构成正当防卫。

\subsubsection*{（一）不法侵害的开始}

从民法上讲，只要不法侵害开始了，哪怕很轻微，就可以进行正当防卫。例如，甲推乙一把，乙反推甲一把。乙属于正当防卫。不过，在刑法上，如前文所述，正当防卫的体系地位是排除违法事由。这就要求正当防卫行为在客观要件上是值得刑罚处罚的危害行为，然后用正当防卫排除掉其可罚的违法性。而乙反推甲一把，这种轻微反击行为在客观要件上不是值得刑罚处罚的危害行为，基于此就不构成犯罪，由此案件就不需要走到“排除违法事由”板块，不需要用刑法的正当防卫来论证其无罪。

基于此，\CJKunderwave{刑法上的正当防卫所制止的不法侵害，要求对法益制造紧迫且较严重的危}\CJKunderwave{险}。例如，甲欲盗窃乙家，在乙家门外踩点考察，此时乙将甲赶走。乙的行为属于民法上的正当防卫，但不属于刑法上的正当防卫，因为乙的赶走行为在客观要件上不属于值得刑罚处罚的危害行为。又如，甲欲盗窃乙家，撬乙家门，乙反击导致甲轻伤。乙的行为先是值得刑罚处罚的危害行为，然后用刑法上的正当防卫来论证无罪。

[注意】设立防卫装置问题。设立防卫装置防卫将来的不法侵害，如果满足以下要求，成立正当防卫：一是防卫手段与不法侵害具有相当性；二是防卫手段不能侵害其他法益，如危害公共安全。由于防卫装置在不法侵害来临时才发挥作用，所以这种防卫在时间上不属于事前防卫。

例1，甲在自家院墙上插了碎玻璃片，小偷乙翻墙时被划伤。甲不属于事前防卫，而属于正当防卫。

例2，甲在自己花房周围私拉电网，电死了小偷。该行为不是正当防卫。因为：第一，该行为危害了公共安全；第二，对盗窃者使用致死手段，不具有相当性。

例3，乙盗窃甲家，甲家看家护院的狗咬住乙。甲不属于事前防卫，而属于正当防卫。

\subsubsection*{（二）不法侵害的结束}

1.结束的内容。不法侵害的结束，不仅要求不法侵害的行为结束，而且要求不法侵害的危险（继续实施侵害的可能性、反扑的可能性）消除，否则，仍认为不法侵害正在进行。（2020年司法解释的规定）

2.不法侵害的危险（反扑可能性）是否消除的判断标准。

% Source PDF physical page 92; see audit/source_coverage.json.
多数说：从行为时、一般人的角度判断。也即，应尊重一般人在当时紧急情况下的认识能力和判断能力，不能要求一般人冷静、理性地判断危险情况。

少数说：从事后专家角度判断。这是事后诸葛亮、上帝视角。

例如（昆山龙哥案，2018年试题“扁担案”的原型），龙哥从宝马车上取刀砍于海明，刀脱手，二人抢刀，于海明抢到刀，为阻挡龙哥，用刀捅龙哥，龙哥被刺中后跑向自己的轿车，于海明追砍龙哥。即使事后专家查明，龙哥跑向轿车时，只是想逃离现场，但是，从一般人的角度看，在当时的情境下，龙哥有可能从车里拿凶器重新投入战斗或开车撞于海明，因此，于海明面临的危险尚未彻底解除。于海明此时追砍龙哥的行为属于防卫适时，而非事后防卫。\footnote{提示：于海明追砍行为虽然属于适时的防卫行为，但是要成立正当防卫，还需要符合限度条件，若超过限度，成立防卫过当。例如，如果在对攻阶段，于海明用刀砍死龙哥，则属于正当防卫。如果在追砍阶段，于海明用刀砍死龙哥，则属于防卫过当。如果无法查明哪个阶段的刀砍导致致命伤，则根据存疑时有利于被告原则，死得越早，对被告人越有利，认定为对攻阶段的刀砍导致致命伤，认定为正当防卫。真实案件中，于海明在追砍阶段，只砍中车身，没有砍中龙哥，表明是对攻阶段的刀砍导致致命伤，因此是正当防卫。}

[注意1]财产犯罪的特殊处理。犯罪人即使取得财物，但在被害人当场还来得及挽回损失的情形下，视为不法侵害尚未结束，可实施正当防卫。例如（2007年第2题），甲使用暴力抢劫到乙的财物后逃跑，乙当场就追，追上后对甲使用暴力夺回财物，致甲重伤。乙的行为属于正当防卫，不属于事后防卫。

[注意2］扭送，是指公民为了防止犯罪人逃避法律追究，有权通过强制手段控制犯罪人，将其送交司法机关。扭送是公民的权利。例如，甲杀人后逃跑。行人乙看到这个场景，抓住甲，将其扭送到公安机关。乙的行为属于扭送，不属于正当防卫、事后防卫，更不构成非法拘禁。（2025年试题）

\subsubsection*{（三）防卫不适时的处理办法}

第一，故意为之，成立故意犯罪（以满足故意犯罪的其他条件为前提）。第二，过失为之，成立过失犯罪（以满足过失犯罪的其他条件为前提）。第三，无故意、过失，成立意外事件。例如，甲使用暴力抢劫乙，被乙打晕在地。乙看到甲昏迷不醒，为了教训甲，又将甲打成重伤。乙后面的行为属于事后防卫，成立故意伤害罪。这种行为不成立防卫过当，不能享受防卫过当的从宽处罚待遇。\footnote{提示：有时候，防卫不适时也会带有假想防卫的特征。如果是故意为之，则认定为防卫不适时，而不认定为假想防卫，因为假想防卫不包括故意犯罪的情形。例如，某晚，甲在家里听到院子里有动静，以为有小偷盗窃，便拎着猎枪出来寻找，在院子里没发现小偷，出门在大街上远远望去有两个人影，认为是小偷逃跑，便开枪打死这两个人。事后查明这两个人是过路行人。此时不存在紧迫的不法侵害，甲竞然朝人开枪，属于故意杀人罪，对甲认定为防卫不适时。虽然甲存在假想，但不认定为假想防卫，因为一旦认定为假想防卫，便无法认定为故意犯罪。}

\subsection*{三、意思条件}

意思条件讨论的是，成立正当防卫，是否要求防卫人主观上具有防卫意思？防卫意思由防卫认识和防卫动机构成。防卫认识（防卫意识），是指防卫人认识到某项合法权利正在遭受不法侵害，认识到自已在制止不法侵害。防卫动机（又称防卫意志），是指防卫人制止不法侵害是出于正当动机，也即为了保护合法权利。

% Source PDF physical page 93; see audit/source_coverage.json.
\subsubsection*{（一）防卫认识}

成立正当防卫，是否要求防卫人具有防卫认识，理论上存在激烈争议，争议主要体现在偶然防卫问题上。这个问题比较疑难，又比较重要，希望勇敢的你耐心学习，多看两遍。

1.偶然防卫的辨认标准：在制止不法侵害这一点上，没有做到主客观相一致，也即客观上制止了不法侵害，但行为人主观上没有认识到这点。偶然防卫的特点是，\CJKunderwave{坏心办}\CJKunderwave{好事}。

2.偶然防卫的类型

（1）故意型偶然防卫。例如，乙欲杀死丙，正要开枪时，被窗外的甲开枪打死。甲没有认识到乙正要杀人，甲只有杀害乙的故意。经事后查明，若甲当时不将乙打死，乙就会将丙打死。（甲\ensuremath{\rightarrow}乙一丙）。甲在客观上制止了乙的不法侵害，但主观上没有认识到这一点。甲属于偶然防卫。甲涉嫌的罪名是故意杀人罪。

（2）过失型偶然防卫。例如，乙欲杀死丙，正要开枪时，被附近的甲开枪打死。甲是猎人，正坐在附近擦枪，不慎走火，打死了乙。经事后查明，若甲当时不将乙打死，乙就会将丙打死。甲在客观上制止了乙的不法侵害，但主观上没有认识到这一点。甲属于偶然防卫。甲涉嫌的罪名是过失致人死亡罪。

3.偶然防卫的观点展示（五星级考点）

（1）防卫认识不要说：成立正当防卫，只要求客观上制止了不法侵害，不要求主观上认识到这一点。因此，无论是故意型偶然防卫，还是过失型偶然防卫，所有的偶然防卫均构成正当防卫。

（2）防卫认识必要说：成立正当防卫，要求主客观相统一，也即不仅要求客观上制止了不法侵害，还要求主观上认识到这一点。因此，所有的偶然防卫均不构成正当防卫。接下来，需要判断偶然防卫行为构成什么犯罪？

防卫认识必要说认为，上述故意型偶然防卫中，甲的行为是坏行为，因此成立故意杀人罪。在此基础上是否构成既遂？既遂结果有两个要求，第一，在事实判断层次，存在死亡结果；第二，在法律评价层次，该死亡结果被评价为危害结果（坏结果）。不能认为，只要是死亡结果就是危害结果。例如，法警执行枪决，制造了死亡结果，但属于好结果。在法律评价上，通过综合比较，甲杀了乙，制止了乙的不法侵害，救了丙，整体上是个好结果，不是危害结果，因此对甲不能定故意杀人罪既遂。剩下只有两个犯罪形态：犯罪中止与未遂。由于甲没有中止行为，根据排除法，对甲只能定故意杀人罪未遂。许多同学对这个结论难以理解，还请多读两遍。\footnote{（1）防卫认识不要说背后的理论依据是结果本位主义（结果无价值论），认为一个行为若制造了好结果，那么这个行为就是好行为。用结果的好坏来评价行为的好坏，以成败论英雄。偶然防卫因为制造了好结果，所以是好行为，应受奖励，认定为正当防卫。（2）防卫认识必要说背后的理论依据是行为本位主义（行为无价值论），认为不能用结果的好坏来评价行为的好坏。行为的好坏与结果的好坏应独立评价。行为的好坏，决定故意犯罪是否成立。结果的好坏，决定故意犯罪是否既遂。故意型偶然防卫中，行为是坏行为，所以成立故意杀人罪，但结果是好结果，所以不构成犯罪既遂，而构成犯罪未遂。简言之，结果本位主义看重结果，行为本位主义看重行为过程。“结果无价值”的意思是“结果是恶的、坏的，是不法结果”，表达的立场是，审查犯罪，要看重不法结果。“行为无价值论”的意思是“行为是恶的、坏的，是不法行为”，表达的立场是，审查犯罪，要看重不法行为。“结果无价值论”也可简称为“结果不法论”，“行为无价值论”也可简称为“行为不法论”。}

防卫认识必要说认为，上述过失型偶然防卫中，甲涉嫌过失致人死亡罪。成立过失犯

% Source PDF physical page 94; see audit/source_coverage.json.
罪，有三项要求，一是有过失行为（坏行为），二是有实害结果（坏结果），三是有因果关系。甲有坏行为，但没有坏结果，因为杀了乙，救了丙。因此，甲不构成过失致人死亡罪，作无罪处理。

[提示】防卫认识不要说和必要说，没有哪个是多数说，二者平分秋色。

\InsertDiagram{figures/manual_p094_1.tex}

4.偶然防卫与假想防卫的比较

相同点：在“制止不法侵害”这一句上，二者都没有做到主客观相一致。区别在于：偶然防卫是客观上制止了不法侵害，主观上没有认识到这一点。假想防卫是客观上不存在不法侵害，主观上以为在制止不法侵害。

\InsertDiagram{figures/manual_table_p094_2.tex}

5.偶然防卫与打击错误的结合

例如（猪队友帮倒忙案），甲、乙共同向丙射击。未料甲击中乙，致乙死亡。事后查明，甲若未致乙死亡，乙便会射中丙。（2016年第52题，2022年试题）

打击错误是三边关系：犯罪人（甲）、预定目标（丙）、实害对象（乙）。甲朝丙开枪，却不慎打死乙。这是打击错误。偶然防卫是双边关系：犯罪人（甲）\ensuremath{\rightarrow}实害对象（乙）。甲打死乙，客观上制止了乙的不法侵害，但甲主观上没有意识到这一点，因此属于偶然防卫。

分析顺序是，先分析打击错误，后分析偶然防卫，因为在犯罪构成体系中，打击错误处在前面的“主观要件”一讲，偶然防卫处在后面的“排除违法事由”一讲。第一，打击错误的角度。（1）具体符合说认为，甲对丙构成故意杀人未遂，对乙构成过失致人死亡，想象竞合，择一重罪论处，定故意杀人罪未遂。（2）法定符合说认为，甲对丙构成故

% Source PDF physical page 95; see audit/source_coverage.json.
意杀人未遂，对乙构成故意杀人既遂（此处可见法定符合说修改案件事实的问题），想象竞合，择一重罪论处，定故意杀人罪既遂。

第二，偶然防卫的角度。偶然防卫分析的是双边关系：犯罪人（甲）\ensuremath{\rightarrow}实害对象（乙），并且以真实的案件事实为分析素材，不考虑法定符合说将甲对乙的过失拟制为故意。（1）防卫认识不要说认为，甲打死乙，构成正当防卫。（2）防卫认识必要说认为，甲打死乙，不构成正当防卫。甲涉嫌过失致人死亡罪。通过综合比较，甲导致乙死亡，救了丙，没有制造坏结果，所以不构成过失致人死亡罪（要构成过失致人死亡罪，要求制造坏的实害结果），对甲作无罪处理。

最后总结结论：（1）甲对丙（预定目标），具体符合说与法定符合说都认为，甲对丙构成故意杀人罪未遂。（2）甲对乙（实害对象），要分两步走。第一步，打击错误。具体符合说和法定符合说都认为，甲对乙构成犯罪，只是罪名不同。第二步，偶然防卫。防卫认识不要说和必要说都认为，甲对乙不构成犯罪，只是理由不同。先立后破，最终应按照偶然防卫的分析结论处理，也即甲对乙不构成犯罪。汇总后的最终结论：甲对丙构成故意杀人罪未遂，甲对乙不构成犯罪。\footnote{读者可以再练习分析此类案件。例1，甲看到乙要向丙射击，为了帮助乙也向丙射击，未料却击中乙，致乙死亡。例2，甲向乙开枪，不慎击中附近的丙，丙死亡，而丙正要枪杀乙，甲对此并不知情。如果甲未将丙打死，丙就会打死乙。}

6.防卫认识的认识错误

结论：防卫人对不法侵害的种类有认识错误，这种认识错误不重要。

例1，甲以为乙在抢劫丙，实际上乙在强奸丙，甲暴力反击乙，甲构成正当防卫。

例2，甲以为乙只是非法侵入住宅，实际上乙是入户盗窃，甲用菜刀将乙砍成重伤。无论针对非法侵入住宅，还是针对入户盗窃，砍成重伤，都属于防卫过当。（2026年试题）

\subsubsection*{（二）防卫动机}

例如，甲看到仇人乙在强制猥亵妇女，出于报复心理攻击乙，将乙打倒。第一，甲认识到了乙在实施不法侵害，甲具有防卫认识，认识到自己在制止不法侵害。第二，甲不是为了保护合法权利而攻击乙，而是出于报复心理。甲不具有防卫动机。

1.成立正当防卫不要求具有防卫动机。上例中，甲虽然缺乏防卫动机，但是甲在客观上制止了乙的不法侵害，可成立正当防卫。（2011年第7题）

2.防卫动机和侵害动机可以并存，不是对立排斥关系。也即，正当动机和不法动机可以并存。例如，甲看到仇人乙在强制猥亵自已的女友，一方面出于报复心理，另一方面出于救助心理，攻击乙，将乙打倒。甲既有报复心理（侵害动机），也有防卫动机，二者并存。甲成立正当防卫。

\subsubsection*{（三）缺乏防卫意思的另两种情形}

1.防卫挑拨

这是指甲欲侵害乙，故意先挑乙（例如辱骂乙、推乙、扇乙耳光），促使乙侵害自己（例如乙用拳头反击甲），然后以正当防卫为借口加害乙。

[结论］甲不成立正当防卫，属于故意犯罪。理由：第一，甲没有防卫意思，只有犯罪故意，属于恶意利用正当防卫制度。第二，甲的挑行为属于不法侵害，乙的反击属于

% Source PDF physical page 96; see audit/source_coverage.json.
正当防卫，对乙的正当防卫不能正当防卫。

例如，甲和朋友乙一同吃夜宵，聊天中，甲戏称乙是肥猪200斤，乙气不过，用身体压制瘦小的甲，致甲呼吸困难、无法动弹，有生命危险。甲用酒瓶砸乙，致乙重伤。甲没有恶意利用正当防卫的意图，不属于防卫挑拨。乙对甲实施严重的不法侵害，甲的反击属于正当防卫。（2023年试题）

2.相互斗殴

这是指双方都出于侵害对方的意图而相互攻击，也即不法侵害对不法侵害。这种情形类似于决斗。例如，甲乙为了争抢女友小芳，双方约架，在大街上相互打架。（2025年试题）

（1）处理结论

[结论一］斗殴无防卫，也即斗殴双方均不成立正当防卫。理由是，正当防卫是正对不正，而相互斗殴是不正对不正，是狗咬狗，都没有防卫意思，都有不法侵害的意思。

[结论二］斗殴有承诺。双方意识到斗殴会使自己受伤，仍自愿参与斗殴，表明对可能出现的伤害结果有预见、有承诺。但是，轻伤害可以承诺放弃，重伤害及生命即使承诺放弃也无效。因此，相互斗殴中，一方将另一方打成轻伤，不构成故意伤害罪，但打成重伤或打死，则成立犯罪。（2010年第7题）

当然，相互斗殴过程中也会出现正当防卫的情形。

例1（一方停止斗殴），甲乙相互斗殴，甲中途停止，求饶或逃跑，乙继续侵害。此时甲可以正当防卫。

例2（一方突然升高级别），甲乙相约赤手空拳斗殴，乙突然掏出手枪欲杀害甲。对于乙的这种严重不法侵害，甲没有承诺，甲此时可以正当防卫。（2026年试题）

（2）相互斗殴与“一方故意伤害、另一方在正当防卫”的区分：

第一，看谁先动手。先动手的一方是不法侵害，还手的一方是防卫行为。

第二，看地点。如果上门打人家，人家的反击是防卫行为。

[提示］不能仅因行为人事先进行防卫准备，就认为其有不法侵害的意图。例如，甲要上门打乙，乙在家里提前准备好菜刀用来防卫。不能因为乙准备了菜刀，就认为乙有不法侵害意图，进而认为乙与甲在相互斗殴。

\subsection*{四、对象条件}

对象条件，是指必须针对不法侵害人本人进行防卫。

\subsubsection*{（一）防卫反击的对象}

1.针对共犯人

例如，甲教唆乙丙丁杀害小芳。乙望风，丙给丁递刀子，丁实施捅刺。小芳能对谁实施防卫反击？

[结论］第一，对直接实行犯可以防卫反击。第二，对幕后的教唆犯、间接正犯不能防卫反击。\footnote{这是因为：第一，教唆行为、间接正犯的指使行为对法益的侵害具有间接性，该类法益侵害缺乏攻击性、紧迫性。第二，正当防卫制止的是正在进行的不法侵害。当直接实行者实施不法侵害，表明教唆行为等已经实施完毕了。}第三，帮助犯如果具有攻击性帮助行为，则对其能够正当防卫。基于此，上述案例中，小芳对甲乙不能防卫反击，对丙丁可以防卫反击。

% Source PDF physical page 97; see audit/source_coverage.json.
2.针对物品

例如，情形一（抢铁锤案）：甲骑摩托车要撞乙，乙用铁锤砸倒摩托车，导致车毁坏，甲摔倒。情形二：甲骑摩托车找到乙，停好车，追打乙，乙逃跑，彻底甩掉甲，发现甲停在路边的摩托车，出于泄愤，端倒摩托车，导致车毁坏。哪种情形，乙属于正当防卫？

[结论］正当防卫的手段，既可以针对不法侵害人的身体，也可以针对不法侵害人的物品。毁坏物品，成立正当防卫的条件是，能够起到制止正在进行的不法侵害行为的效果。基于此，情形一中，乙成立正当防卫。情形二中，乙不构成正当防卫，构成故意毁坏财物罪（以财物价值达到数额较大为条件）。（2022年试题）

\subsubsection*{（二）三种特殊情形}

1.甲故意伤害乙，乙情急之下将丙的珍贵花瓶拿起砸向甲，砸中甲，导致花瓶毁损。

结论：乙对甲构成正当防卫，对丙构成紧急避险（以符合紧急避险其他条件为前提），二者想象竞合，优先认定为正当防卫。虽然乙对丙构成紧急避险，不构成故意毁坏财物罪，但民事上应赔偿丙。（2016年第6题）

2.甲故意伤害乙，乙情急之下将丙的珍贵花瓶拿起抵挡（没有反击），花瓶毁损。

结论：第一，乙对甲不构成正当防卫，因为乙只是抵挡，不是反击甲。第二，乙对丙构成紧急避险，因为乙为了保护自己身体法益，不得已损害丙的财产法益。如果乙不拿丙的花瓶，丙的花瓶不会受损。（2024年试题）

3.甲拿起丙的珍贵花瓶扔向乙，欲伤害乙，乙用手臂抵挡，花瓶掉地上，摔碎。

结论：乙的行为本身没有违法性，因此对甲不成立正当防卫，对丙不成立紧急避险。理由：正当防卫、紧急避险都是排除违法事由，要成立正当防卫、紧急避险，要求行为先要有违法性（法益侵害性）。第一，乙只是单纯抵挡，不是反击甲。第二，花瓶摔碎，这一法益侵害结果应归责于甲，而不应归责于乙。乙的行为本身没有违法性，在客观要件这一步就得出无罪结论。提示，该案与上文“抢铁锤案”区别在于，该案中，乙只是单纯抵挡物品，不是反击甲的不法侵害。而“抢铁锤案”中，乙用铁锤砸摩托车，是一种反击行为，反击甲的进攻。

\subsection*{五、限度条件}

第20条第2款正当防卫明显超过\CJKunderwave{必要限度}造成\CJKunderwave{重大损害}的，应当负刑事责任，但是应当减轻或者免除处罚。

\subsubsection*{（一）判断标准}

1.第一位判断标准：必要性标准。这是指，只要是制止不法侵害的必要手段，就不过当。如果不必要，就有可能过当。如何判断防卫手段是否具有必要性？

多数说（司法解释）：根据行为时、一般人的角度判断。也即，在紧急时刻，不应当苟求防卫人像黄飞鸿一样能够冷静判断，点到为止。

少数说：根据事后专家角度判断。这是上帝视角、事后诸葛亮。

例如，甲持枪抢劫银行，警方警示无效，情况紧急，被迫开枪击毙甲。事后专家查明，甲的枪是毫无杀伤力的仿真枪。但是，从行为时、一般人的角度看，无法知道甲的枪是仿真枪，因此，警察开枪是必要手段。

% Source PDF physical page 98; see audit/source_coverage.json.
2.第二位判断标准：相当性标准。这是指，根据比例原则，不法侵害越严重，则防卫级别可以越高。不法侵害越轻，则防卫级别也应越低。

例1，狗蛋屡次半夜闯入村里寡妇小芳家，使用暴力强奸小芳。有一次，狗蛋又半夜闯入小芳家，要求小芳脱衣服。小芳孤立无援，欺骗狗蛋喝杯枸杞茶（实际是毒药），狗蛋不知情中毒身亡。如果将相当性作为第一位标准，则会认为，狗蛋尚未实施暴力，小芳就将其毒死，应成立防卫过当。而将必要性作为第一位标准，则会认为，此时的小芳别无他法，只有毒死狗蛋才能制止其不法侵害。若等到狗蛋实施暴力，则小芳便失去制止不法侵害的可能性了。

例2，小偷翻窗，下半身卡在外面。主人用菜刀将小偷砍成重伤。主人构成防卫过当。(2019年试题）

\subsubsection*{（二）防卫过当}

法条规定：防卫过当，应当负刑事责任，但应当减轻或者免除处罚。

\InsertDiagram{figures/manual_table_p098_1.tex}
\SourceFootnote{\textbf{表注①}\quad 对“防卫过当”，也可以阶段化理解，有客观意义上的防卫过当，也有最终意义上的防卫过当。这一点，与“犯罪”概念类似。}

\subsection*{六、最厉害的正当防卫（特殊正当防卫）}

第20条第3款对正在进行行凶、杀人、抢劫、强奸、绑架以及其他严重危及人身安全的暴力犯罪，采取防卫行为，造成不法侵害人伤亡的，不属于防卫过当，不负刑事责任。

\subsubsection*{（一）条文属性}

该款规定不是特殊规定，而是正常规定，并没有改变正当防卫的一般成立条件。要成立最厉害的正当防卫，也要符合一般正当防卫的成立条件。例如，要符合时间上的适时性、手段上的必要性和相当性等。之所以造成死亡也不过当，是因为面临的不法侵害是严重危及人身安全的暴力犯罪，所以防卫的程度级别相应提高，造成死亡也是必要的、相当的。

\subsubsection*{（二）具体解释}

1.行凶。这是指严重危及人身安全的暴力行为。成立行凶，不要求携带凶器。

% Source PDF physical page 99; see audit/source_coverage.json.
2.“杀人、抢劫、强奸、绑架”。（1）这四项是指犯罪行为，而不是具体罪名。例如，“抢劫”包括抢劫枪支、弹药，“绑架”包括以绑架手段拐卖妇女、儿童。（2）这四种犯罪行为必须是暴力手段严重危及人身安全，否则只能实施一般的正当防卫。例如，母亲故意不喂养婴儿或投毒杀人，对母亲、投毒者不能实施最厉害的正当防卫，致其重伤或死亡。若致其重伤、死亡，则属于防卫过当。但是可以实施一般的正当防卫。又如，用麻醉方法抢劫，携带凶器抢夺拟制的抢劫（只对物暴力），对抢劫者不能实施最厉害的正当防卫，致其重伤、死亡。（2021年试题）

3.“其他严重危及人身安全的暴力犯罪”，应当是与杀人、抢劫、强奸、绑架行为相当，并具有致人重伤或者死亡的紧迫危险和现实可能的暴力犯罪。例如，劫持航空器罪、放火罪、爆炸罪、暴动越狱罪，对这些犯罪可以进行最厉害的正当防卫。\footnote{以上内容参见2020年司法解释：2020年8月28日最高人民法院、最高人民检察院、公安部《关于依法适用正当防卫制度的指导意见》。}

[大总结]追小偷的问题。

1.追小偷的行为包括三种类型。

（1）追财行为。这是指，追小偷是为了夺回财物。这又分为两种情形：①正在追的途中，小偷自已摔死。正在追的行为本身合法，无罪，不需要用正当防卫来论证无罪。②追上，夺回财物。这是正当防卫。

（2）扭送行为。这是指，追小偷是为了扭送。这是合法行为，但条件是，只能以控制小偷的人身自由为目的和限度，不能夹带“私货”，如泄愤伤害等。如果题中没有特别交代，追小偷的行为一般指扭送行为。

（3）不法侵害行为。这是指，追小偷是为了伤害侮辱等，或者先是为了扭送，追上控制后，又实施伤害侮辱。这都属于不法侵害行为。（2020年试题“偷狗案”）

2.考察因果关系。

（1）小偷逃跑过程中，不慎摔死。

①如果追小偷的行为属于追财行为、扭送行为等合法行为，则与小偷死亡没有因果关系。这是因为，因果关系的“因”必须是危害行为。（2010年第3题）

②如果追小偷的行为属于不法侵害行为，则根据“介入因素两步走”标准判断，看介入因素（被害人自身行为也即小偷摔倒）是否异常。对此需要情景化判断，看追的行为（不法侵害行为）的紧迫性和严重性。如果甲拿着棍棒要殴打小偷，紧追不舍，小偷逃跑时摔倒，便不异常，那么死亡结果与追的行为具有因果关系。

（2）追捕者追的过程中，不慎摔死，或被车撞死。

结论是，追捕者的死亡与小偷没有因果关系。理由是，小偷的逃跑行为没有给追捕者制造危险，不是危害行为，不符合“因”的要求。在此前提下，不需要判断介入因素是否异常。（2013年主观题，2019年试题）

引申情形：甲追杀乙，紧追不舍，乙逃跑，闯红灯，被车撞死。介入因素（车祸）不异常，因此甲的行为与乙的死亡有因果关系。

3.考察作为义务。

小偷逃跑过程中，跳河，有淹死危险。追捕者有无救助义务？

% Source PDF physical page 100; see audit/source_coverage.json.
判断标准：追赶行为对小偷是否制造紧迫危险，小偷是不是被迫跳河。如果是，则有救助义务。如果不是，例如追捕者距离小偷还有较远距离，小偷此时跳河，属于被害人自陷风险，自己负责，追捕者没有救助义务。注意，此时不需要考察追的行为是否合法，因为多数说认为，即使是合法行为，若对法益制造紧迫危险，也会成为先行行为而产生作为义务。

\section{紧急避险}

[问题］影片《泰坦尼克号》的片尾，杰克和罗丝在海上共抓一块木板，木板只能浮住一个人。假如杰克为了救自己，一拳将罗丝打下去。对杰克该如何处理？（海上木板案）

第21条第1款为了使国家、公共利益、本人或者他人的人身、财产和其他权利免受正在发生的危险，不得已采取的紧急避险行为，造成损害的，不负刑事责任。

\subsection*{一、起因条件}

1.危险的来源

紧急避险的起因条件是必须存在现实危险。人面临的危险主要分为两种。

（1）自然灾害、野生动物袭击、疾病。例如，野狗追咬狗蛋，狗蛋迫不得已闯入小芳家。狗蛋属于紧急避险，不构成非法侵入住宅罪。

（2）他人的不法侵害行为。例如，狗剩追砍狗蛋，狗蛋迫不得已闯入小芳家。狗蛋属于紧急避险。\footnote{紧急避险的“危险”还包括被害人自陷的危险。例如，甲关起门，在家里自杀。乙为了抢救甲，被迫砸坏门窗（造成较大财产损失），破门而入抢救甲。乙的行为属于紧急避险。避免危险的主要方式是攻击无辜第三人，例如闯进小芳家。这种方式被称为攻击型紧急避险。此外，避免危险的方式还有反击危险源。例1，野狗追咬狗蛋，狗蛋反击野狗。这种方式被称为防御型紧急避险。例2，狗剩追砍狗蛋，狗蛋反击狗剩。这属于防御型紧急避险。这种防御型紧急避险避免的危险是人的不法侵害，是在制止人的不法侵害，所以也被称为正当防卫。可见，正当防卫是防御型紧急避险中的一种情形。}

[注意1]对他人合法行为制造的危险不能紧急避险。例如，警察抓捕小偷，小偷为了躲避警察抓捕，夺过第三人的摩托车骑走或闯入第三人的住宅，不构成紧急避险。（2016年第6题）

[注意2］负有特定职责的人在面临自己职务、业务带来的危险时，不能进行紧急避险。例如，消防队员面对火灾时，不能进行紧急避险。警察面对自己职责范围内的犯罪时，不能进行紧急避险。

[注意3]可以为了保护国家法益、公共法益而紧急避险。例1，消防员为了消灭大楼火灾，不得已闯入他人院子使用水管，属于紧急避险。灭火是保护公共法益。例2，警察为了抓捕小偷，不得已踩踏村民的麦田，属于紧急避险。抓小偷是保护社会法益。（2024年试题）

2.危险的现实性

客观上不存在危险，行为人误以为存在危险，进行避险行为，属于假想避险。与假想防卫的处理办法相同，对于假想避险，如果存在过失，就是过失犯罪；如果没有过失，就是意外事件。例如，狗蛋在大街上不慎踩到狗屎，斜靠在电线杆上用力甩狗屎。狗剩以为狗蛋触电了，用木棒打倒狗蛋。狗剩属于假想避险。

% Source PDF physical page 101; see audit/source_coverage.json.
\subsection*{二、时间条件}

时间条件是危险正在发生。所谓正在发生，是指法益受到的危险已经产生但尚未消除。如果事前或事后避险，属于避险不适时，处理办法与防卫不适时相同。\footnote{紧急避险要求“危险正在发生”，正当防卫要求“不法侵害正在进行”。后者的紧迫程度比前者要更紧迫一些，前者可以更轻缓、更遥远。例如，古代，山顶客栈里，店小二甲偷听到几个客人在密谋“吃完饭，下山，杀死山脚下的乙”。甲不得已用蒙汗药弄晕客人，然后下山告知乙赶紧逃离。甲弄晕客人，是在制止不法侵害，但该不法侵害不够紧迫，所以甲不属于正当防卫，而属于事前防卫。但是，该不法侵害也是一种危险，该危险正在发生，甲弄晕客人，属于防御型紧急避险。}

\subsection*{三、意思条件}

意思条件讨论的是，成立紧急避险，是否要求避险人主观上具有避险意思？避险意思由避险认识和避险意图构成。避险认识（避险意识），是指避险人认识到某项合法权利正在遭受危险，认识到自己在避免危险。避险意图（又称避险意志、避险动机），是指避险人实施避险行为是出于正当意图、动机，也即为了保护合法权利。

\subsubsection*{（一）避险认识}

成立紧急避险，是否要求避险人具有避险认识，理论上存在激烈争议，争议主要体现在偶然避险问题上。

\InsertDiagram{figures/manual_table_p101_1.tex}

\subsubsection*{（二）避险意图（避险动机）}

例如，甲看到一条野狗在追咬乙，看到丙（与甲有过节）捧着花瓶，情况紧急，出于报复丙的心理，夺过花瓶砸向野狗，救了乙，但导致花瓶打碎。第一，甲认识到了乙面临危险，甲具有避险认识，认识到自已在避免危险。第二，甲实施避险行为，不是出于保护合法权利的动机，而是出于报复心理。甲不具有避险意图。

[结论］成立紧急避险不要求具有避险意图。上例中，甲虽然缺乏避险意图，但是甲在客观上避免了乙遭受危险，可成立紧急避险。

[提示】避险意图和侵害意图可以并存，不是对立排斥关系。也即，合法动机和不法动机可以并存。例如上例中，甲可以一方面出于报复丙的心理，另一方面出于救助乙的心理，用花瓶砸野狗。甲既有报复心理，也有避险意图，二者并存。甲成立紧急避险。（2015年第4题）

% Source PDF physical page 102; see audit/source_coverage.json.
\subsection*{四、“不得已”条件}

这是指避险手段只能是不得已而为之。不得已，意味着没有其他合理的方法可以选择。“不得已”条件也称为补充性条件，意思是避险手段只能是补充手段，而不能是优先手段。这一点与正当防卫不同，正当防卫中，面临不法侵害，原则上，防卫人没有躲避义务，防卫反击是优先手段，不需要不得已而为之。

[注意】受强制的紧急避险。这是指受他人强制实施紧急避险的情形。

例如，甲看到乙带着儿子在银行附近散步，上前持刀绑架了乙的孩子，要求乙立刻去抢劫附近银行，否则杀害其子。乙为了挽救儿子的生命而抢劫了银行。只要这种受强制的行为符合紧急避险的其他条件，便可以成立紧急避险，也即不构成胁从犯。乙的行为成立紧急避险。

\subsection*{五、限度条件}

第21条第2款紧急避险超过必要限度造成不应有的损害的，应当负刑事责任，但是应当减轻或者免除处罚。

要求：保护的法益\ensuremath{\ge}损害的法益

1.法益的种类上，大体而言：生命法益\ensuremath{>}身体健康法益\ensuremath{>}人身自由法益\ensuremath{>}财产法益。

法益衡量上，不存在国家法益\ensuremath{>}公共法益（社会法益）\ensuremath{>}个人法益，也不存在个人法益\ensuremath{>}社会法益\ensuremath{>}国家法益。例如，为了保护个人法益，可以通过损害公共法益或国家法益来紧急避险。又如，为了保护国家法益，可以通过损害某些个人法益来紧急避险。（2016年第6题)

2.财产法益可以画等号。也即，为了保护某个财物，牺牲另一同价值的财物，成立紧急避险，而非避险过当。（2015年第4题）

3.生命法益不能画等号。也即，不能为了保护一个人的生命，牺牲另一个人的生命。

例如，在本节开头“海上木板案”中，由于双方获救机会均等，杰克不构成紧急避险，而构成故意杀人罪。

问题：能否牺牲一个人的生命以保护其他多数人的生命？原则上不能，因为人的生命无价，不能成为被利用的手段。例外，当这个人注定被牺牲，也即牺牲地位被特定化了，可以对其实施紧急避险（多数观点）。例如，“9·11”事件中，恐怖分子劫持客机欲撞向地面大楼。客机上的乘客注定会被牺牲，不撞向地面大楼，也会死，可以用导弹击落，以避免地面大量人员受伤害。

[总结1]正当防卫与紧急避险的比较。

\InsertDiagram{figures/manual_table_p102_1.tex}

% Source PDF physical page 103; see audit/source_coverage.json.
\ContinuedTable{figures/manual_table_p103_1.tex}

[总结2]正当防卫与紧急避险的内部关系。

\InsertDiagram{figures/manual_table_p103_2.tex}

[注意］对正当防卫能否紧急避险？原则上不可以。（2024年试题）

例如，甲杀乙，乙防卫反击、反杀，甲为了躲避反杀，抢来旁边丙的摩托车，骑走。甲的抢车行为不构成紧急避险，构成抢劫罪。这是因为，甲面临的危险（乙的防卫反击）是甲故意自已招来的，刑法对其不予保护。如果乙的反杀变成事后防卫，则甲既可以紧急避险，也可以正当防卫。

\section{被害人承诺}

[问题］狗蛋涉嫌犯罪被抓，狗剩欺骗狗蛋的妻子小芳：“如果愿意陪我一夜，我保证将你老公捞出来！”小芳信以为真，救夫心切，便答应照办。事后，狗剩玩失踪。狗剩是否构成强奸罪？（欺骗上床案）

% Source PDF physical page 104; see audit/source_coverage.json.
\subsection*{一、现实的被害人承诺}

现实的被害人承诺，是指被害人同意他人对其加害，基于此，他人不构成犯罪。例如，甲同意乙毁坏自己的财物，乙的毁坏行为便不构成犯罪。被害人承诺的成立要件：

\subsubsection*{（一）承诺的权限与范围}

被害人对承诺的\CJKunderwave{法益}具有处分权限。例如，甲对乙说：“那辆车不是我的，但我同意你把它砸毁。”乙若毁坏该车便构成犯罪。

被害人对承诺的法益具有一定限度，超出这个限度，即使有被害人的承诺，行为人也有罪。

1.财产可以承诺放弃。例如，甲承诺让乙毁坏自己的车辆，乙便毁坏了车辆，乙无罪。

2.名誉可以承诺放弃。例如，甲承诺让乙诽谤自己，乙便诽谤甲，乙无罪。

3.自由可以承诺放弃。例如，甲让乙拘禁自己，乙拘禁了甲，乙无罪。

4.身体权在轻伤的范围内可以放弃，超出轻伤范围，不可放弃。

例1，甲承诺让乙打自已两耳光，乙便打了甲两耳光，乙无罪。如果甲承诺让乙砍掉自己的手臂，乙便砍掉了甲的手臂（重伤），乙构成故意伤害罪。

例2，甲欲加入黑社会团伙，团伙老大要求测试甲的抗打击能力。甲同意，老大殴打甲，致其重伤。老大构成故意伤害罪。（2024年试题）

[注意1]《人体损伤程度鉴定标准》认为砍小手指属于轻伤，可以承诺。（2024年试题）

[注意2］如果承诺重伤害，是为了保护另一重大法益，则承诺有效。例如，甲承诺摘取自己的肝，用来挽救其他病人生命。

5.生命不可以承诺放弃。例如，甲承诺让乙杀了自己，乙便杀了甲，乙构成故意杀人罪。

[注意1]超出承诺范围。经承诺实施的行为不得超出承诺预设的范围。例如，甲仅承诺让乙打一耳光，而乙将甲打成脑震荡，乙便构成故意伤害罪。

[注意2］虽然经被害人承诺，不构成犯罪，但是如果触犯其他犯罪，则以其他犯罪论处。例如，甲同意乙烧毁自已的房屋，虽然乙不构成故意毁坏财物罪，但是如果危害公共安全，会构成放火罪。

\subsubsection*{（二）承诺的时间}

被害人必须有现实的承诺，并且必须事前作出，事后承诺于事无补。例如，甲盗窃了乙的财物。乙事后发现甲是自已的侄子，向警方表示自己愿意将财物给甲。但是这种事后承诺不影响甲盗窃罪的成立。

\subsubsection*{（三）承诺的能力}

被害人对承诺事项的意义、范围有理解能力。基于此，幼儿、精神病患者的承诺无效。例如，幼女同意与甲男发生性行为。甲男仍构成强奸罪。

\subsubsection*{（四）承诺的意思表示}

被害人的承诺必须是其真实意思表示。受到欺骗、胁迫或乘人之危而作出的承诺无效。例如，甲胁迫乙女：“若不跟我上床，就将你的裸照曝光。”乙被迫答应。乙的承诺无效，甲构成强奸罪。比较疑难的是欺骗的情形。

% Source PDF physical page 105; see audit/source_coverage.json.
1.甲欺骗乙，导致乙产生认识错误

甲欺骗乙，导致乙产生认识错误，乙基于认识错误而作出承诺，承诺放弃法益。乙的承诺是否有效？判断标准：

（1）事实认识错误。甲虚构事实或隐瞒真相（事实），导致乙针对事实产生认识错误，乙基于事实认识错误而作出的承诺无效。这里的事实是指与法益相关联的事实。

例1，甲冒充乙的丈夫欲和乙发生性关系，乙误以为是自己丈夫便同意。乙对重要的法益事实（行为对象的身份）发生认识错误，其同意发生性关系的承诺是无效的。甲构成强奸罪。

例2，甲欺骗乙：“你家小狗患了狂犬病，赶紧交给我来处理掉。”乙相信照办。甲将狗杀死。乙对法益关联事实有认识错误，承诺无效。甲构成故意毁坏财物罪。

例3（2014年第15题），医生甲欺骗乙说：“你的孩子丙需要移植眼角膜，你愿意为你的孩子捐献眼角膜吗？”乙答应。甲却将乙的眼角膜移植给丁。乙对法益关联事实（自已眼角膜的用途）有认识错误，承诺无效。甲构成故意伤害罪。当然，乙的动机没有实现，也有动机错误。事实认识错误与动机错误不是对立排斥关系，可以并存，此时应以事实认识错误的处理结论为准。

（2）动机错误。甲在事实层面没有欺骗乙，乙对与法益关联的事实没有认识错误，但乙的背后动机没有实现，存在动机错误。动机错误不影响承诺的有效性。

例如，甲（男导演）欺骗乙（女演员）：“陪我睡觉，就给女一号。”乙答应。后甲并没有让乙演女一号。乙对法益关联事实（与谁发生性关系）没有认识错误，只是动机落空，存在动机错误。乙的承诺有效，甲不构成强奸罪。同理，本节开头“欺骗上床案”中，狗剩不构成强奸罪。

[提示］事实认识错误与动机错误这种区分标准，是学界的通说，也是法考的官方观点。

2.乙自己产生认识错误

（1）甲不知道乙自己产生认识错误，乙的承诺对甲而言是有效的。

例如（2019年试题），甲用短信问邻居乙，可否砍伐乙后院的一棵大树。乙用短信本想回复不行，不小心漏写“不”，发成“行”。甲便砍伐了乙的树。乙的承诺对甲而言是有效的。

（2）甲知道乙自己产生了认识错误。在此前提下，如果甲实施欺骗行为（作为的欺骗或不作为的欺骗），使乙作出错误承诺或者强化错误承诺，则乙的承诺无效。如果甲没有实施欺骗行为，只是单纯利用乙的错误，则乙的承诺有效。

例如（2019年试题），乙养的马生病了，乙认为无药可救，请求兽医甲实施安乐死。甲知道有特效药可以救治，故意不告知，实施了安乐死。甲作为兽医，基于职业要求，有告知真相义务，故意不告知，属于不作为的欺骗，因此乙的承诺无效。甲构成故意毁坏财物罪。\footnote{以上内容参见王钢：《动机错误下的承诺有效性问题研究》，载《中外法学》2020年第1期。王钢老师是清华大学法学院博导，是法考官方教材的副主编。所以，这一块知识点的观点，以王钢老师的观点为准。}

% Source PDF physical page 106; see audit/source_coverage.json.
\InsertDiagram{figures/manual_p106_1.tex}

\subsection*{二、推定的被害人承诺}

推定的被害人承诺，是指现实中没有被害人的承诺，但是推定被害人得知真相后会作出承诺，基于这种推定的承诺作出的行为不构成犯罪。例如，邻居不在家，屋内发生火灾，甲为了抢救，闯入屋内。推定的被害人承诺的成立条件：

1.被害人没有现实的承诺。例如，店员暂时离开商店，顾客按照标明的价格将现金放在柜台，拿走商品。店员有现实的承诺，不过是默示的承诺。这属于前文的被害人承诺，而不属于推定的被害人承诺。

2.推定被害人得知真相后会承诺。这种推定以一般人的合理意愿为标准，而不以被害人的实际意愿为标准。

例1，甲看到乙女落水，将乙女救上来，救助中造成乙轻伤。乙女扇甲一耳光：“我失恋了，要自杀，为什么救我？”甲的行为无罪。

例2，甲受伤昏迷，急需输血。医生输血将甲抢救过来，但甲责问医生：“依我的宗教信仰，不能输入别人的血，我宁愿死。”医生的行为无罪。

3.必须是为了被害人的一部分法益牺牲其另一部分法益，但所牺牲的法益不得大于所保护的法益。

4.行为所指向的法益必须是被害人有处分权的法益。

[注意]推定的被害人承诺可以与紧急避险产生竞合。也即，一个行为既属于推定的被害人承诺，也属于紧急避险。例如，没有亲属在场、陷入昏迷的甲急需截肢，否则有生命危险。医生实施了截肢。这既属于紧急避险，也属于推定的被害人承诺。
